%%%%%%%%%%%%%%%%%%%%%%%%%%%%%%%%%%%%%%%%%
% Lachaise Assignment
% LaTeX Template
% Version 1.0 (26/6/2018)
%
% This template originates from:
% http://www.LaTeXTemplates.com
%
% Authors:
% Marion Lachaise & François Févotte
% Vel (vel@LaTeXTemplates.com)
%
% License:
% CC BY-NC-SA 3.0 (http://creativecommons.org/licenses/by-nc-sa/3.0/)
%
%%%%%%%%%%%%%%%%%%%%%%%%%%%%%%%%%%%%%%%%%

%----------------------------------------------------------------------------------------
%	PACKAGES AND OTHER DOCUMENT CONFIGURATIONS
%----------------------------------------------------------------------------------------

\documentclass{article}

\usepackage[margin=1cm, bottom=2cm]{geometry}

% \input{structure.tex} % Include the file specifying the document structure and custom commands

% \usepackage{indentfirst}\setlength{\parindent}{0em}
\usepackage{listings}
\usepackage{float}
\usepackage{amsmath,amsfonts,stmaryrd,amssymb} % Math packages
\usepackage{graphicx}
\usepackage{subcaption}
\usepackage{tikz}

% add linking w/ blue color
\usepackage{xcolor}
\usepackage{hyperref}

\hypersetup{
    colorlinks=true,
    linkcolor=blue,
    filecolor=magenta,
    urlcolor=blue,
}

\numberwithin{figure}{section}

\title{Draft}

\begin{document}
\maketitle

\section{Josephson Junction}

\subsection{Newton Update}
The nonlinear system of equations to be solved is
\begin{align}
    \begin{split}
        -\alpha_{x1}E_{x,i+\tfrac12,j-1,k}^{n+\frac12}
        + \beta_{x1}E_{x,i+\tfrac12,j,k}^{n+\frac12}
        - \gamma_{x1}E_{x,i+\tfrac12,j+1,k}^{n+\frac12}
        +\frac12 J_{\mathrm{J},x,i+\tfrac12,j,k}^{n+\frac12}
         & = p_{x1}E_{x,i+\tfrac12,j,k}^{n}                                                                       \\
         & \quad +\frac{1}{\Delta y}\Big(H_{z,i+\tfrac12,j+\tfrac12,k}^{n}-H_{z,i+\tfrac12,j-\tfrac12,k}^{n}\Big) \\
         & \quad -\frac{1}{\Delta z}\Big(H_{y,i+\tfrac12,j,k+\tfrac12}^{n}-H_{y,i+\tfrac12,j,k-\tfrac12}^{n}\Big) \\
         & \quad +\frac{q_{x1}}{\Delta x\Delta y}\Big(E_{y,i+1,j-\tfrac12,k}^{n}-E_{y,i,j-\tfrac12,k}^{n}\Big)    \\
         & \quad -\frac{r_{x1}}{\Delta x\Delta y}\Big(E_{y,i+1,j+\tfrac12,k}^{n}-E_{y,i,j+\tfrac12,k}^{n}\Big)    \\
         & \quad -\frac12 J_{\mathrm{J},x,i+\tfrac12,j,k}^{n},
    \end{split}
    \label{eq:first-e-ex-adi-tridiagonal-josephson-newton}
\end{align}
where
\begin{equation}\label{eq:josephson-new-current}
    J_{\mathrm{J},x}^{n+\frac12}=j_{c,x}\sin\!\left[\phi_x^n+\frac{\Delta t}{4}\eta_x\left(E_x^{n+\frac12}+E_x^n\right)\right].
\end{equation}

Formulate into matrix form, along the \(y\)-line at fixed \(i,k\), as
\begin{equation}\label{eq:josephson-matrix}
    \mathbf{A}\mathbf{E}+\mathbf{J}(\mathbf{E})=\mathbf{b}.
\end{equation}
The unknown is the new electric field on that line,
\begin{equation}\label{eq:josephson-unknown}
    E_j=E_{x,i+\tfrac12,j,k}^{n+\frac12}.
\end{equation}
\(\mathbf{A}\) is the tridiagonal matrix (or low-rank perturbed tridiagonal matrix if a periodic boundary condition is used) formed from the linear ADI stencil, and \(\mathbf{J}(\mathbf{E})\) is a diagonal matrix, 
\begin{equation}
    \mathbf{J}(\mathbf{E}) = \operatorname{diag}\!\left(\mathbf{J}_1(\mathbf{E}),\ldots,\mathbf{J}_M(\mathbf{E})\right),
\end{equation}
where \(M\) is the number of unknowns on the line, and its \(j\)-th diagonal element is given by
\begin{equation}\label{eq:josephson-J-of-E}
    \mathbf{J}_j(\mathbf{E})=\frac{j_{c,x}}{2}\sin\!\left[\phi_{x,i+\tfrac12,j,k}^{n}+\frac{\Delta t}{4}\eta_x\left(E_j+E_{x,i+\tfrac12,j,k}^{n}\right)\right],
\end{equation}
so \(\mathbf{J}(\mathbf{E})\) changes only the diagonal. $\mathbf{b}$ is simply the right-hand side.

Define the equation to be solved as
\begin{equation}\label{eq:josephson-residual}
    \mathbf{F}(\mathbf{E})=\mathbf{A}\mathbf{E}+\mathbf{J}(\mathbf{E})-\mathbf{b}=\mathbf{0}.
\end{equation}
The Jacobian matrix of the nonlinear term is
\begin{equation}\label{eq:josephson-current-jacobian}
    J_{\mathbf{J}}(\mathbf{E}) = \frac{\partial \mathbf{J}(\mathbf{E})}{\partial \mathbf{E}} = \operatorname{diag}\!\left(d_1(\mathbf{E}),\ldots,d_M(\mathbf{E})\right),
\end{equation}
where
\begin{equation}\label{eq:josephson-current-jacobian-diagonal}
    d_j(\mathbf{E})= \frac{\partial \mathbf{J}_j(\mathbf{E})}{\partial E_j} = \frac{\Delta t}{8}\eta_x j_{c,x}\cos\!\left[\phi_{x,i+\tfrac12,j,k}^{n}+\frac{\Delta t}{4}\eta_x\left(E_j+E_{x,i+\tfrac12,j,k}^{n}\right)\right].
\end{equation}
Here \(\eta_x\) and \(j_{c,x}\) are evaluated at edge \(i+\tfrac12,j,k\). The Jacobian of the full residual is therefore
\begin{equation}\label{eq:josephson-residual-jacobian}
    J_{\mathbf{F}}(\mathbf{E})=\frac{\partial\mathbf{F}(\mathbf{E})}{\partial\mathbf{E}}=\mathbf{A}+J_{\mathbf{J}}(\mathbf{E}),
\end{equation}
which retains the tridiagonal structure. Thus, the Newton update is
\begin{equation}
    \mathbf{E}^{(k+1)} = \mathbf{E}^{(k)} + \Delta\mathbf{E}^{(k)},
\end{equation}
where
\begin{equation}
    \Delta\mathbf{E}^{(k)} = -\left(\mathbf{A}+J_{\mathbf{J}}(\mathbf{E}^{(k)})\right)^{-1}\mathbf{F}(\mathbf{E}^{(k)}).
\end{equation}
In practice, we solve the linear system
\begin{equation}
    \left(\mathbf{A}+J_{\mathbf{J}}(\mathbf{E}^{(k)})\right)\Delta\mathbf{E}^{(k)} = -\mathbf{F}(\mathbf{E}^{(k)}).
\end{equation}

\end{document}